% =============================================================================
% fig4-7-dma-subsystem.tex --- 图 4.7 DMA 子系统结构 (§4.10)
%
% v5.1 修订: 与 fig3-3 差异化:
%   - fig3-3 = 对外接口黑盒视角 (DMA 子系统 1 master + 1 slave)
%   - fig4-7 = 内部 cmd / data 双通道展开视角
%               * axi_dm IP 展开为 MM2S cmd FIFO + MM2S data
%                                 + S2MM cmd FIFO + S2MM data + sts FIFO
%               * mm2s_arb 居 axi_dm 上方 (3 ctrls 汇聚)
%               * axi_m_mux 3-to-1 聚合 + DFE 旁路 (旁路从最右侧绕)
%   - 节点文字一两个词, 通道方向 / 仲裁策略放节点旁小字标注 (无括号)
%
% v5.2 字体优化: 全局放大字号; sts 双输入分线; sts→ctrls 标签移出 IP shell
% =============================================================================
\documentclass[border=4pt]{standalone}
% =============================================================================
% _tikz-style.tex --- FLUX_CNN 论文 TikZ 共享样式包
%
% 用法（每张 figX-Y.tex 头部 \input 本文件）:
%   \documentclass[border=4pt,varwidth=18cm]{standalone}
%   % =============================================================================
% _tikz-style.tex --- FLUX_CNN 论文 TikZ 共享样式包
%
% 用法（每张 figX-Y.tex 头部 \input 本文件）:
%   \documentclass[border=4pt,varwidth=18cm]{standalone}
%   % =============================================================================
% _tikz-style.tex --- FLUX_CNN 论文 TikZ 共享样式包
%
% 用法（每张 figX-Y.tex 头部 \input 本文件）:
%   \documentclass[border=4pt,varwidth=18cm]{standalone}
%   \input{_tikz-style.tex}
%   \begin{document}
%   \begin{tikzpicture}[<额外局部样式>]
%     ...
%   \end{tikzpicture}
%   \end{document}
%
% 风格基线: IEEE 期刊配色 / 白底黑字 / Times New Roman + Microsoft YaHei /
%           禁卡通 / 禁渐变 / 禁阴影 / 禁 3D / 禁 emoji
% 与 figures-rendered/_style.py (matplotlib) 配色一致, 保证全文图风格统一.
% =============================================================================

% ---- 字体（XeLaTeX 必须）----
% 中文字体：Microsoft YaHei（保持原配置，本次只调整英文）.
% 英文字体：Times New Roman, 把 \sffamily / \rmfamily 同时映射到 Times,
%           保证 TikZ 节点无论用哪个 family, 拉丁字符都是 Times New Roman.
% 数学公式（mathmode）保留 latin modern math, 与 Times 视觉差异可接受.
\usepackage{amsmath}
\usepackage{amssymb}
\usepackage{xeCJK}
\setCJKmainfont{Microsoft YaHei}
\setCJKsansfont{Microsoft YaHei}
\setmainfont{Times New Roman}
\setsansfont{Times New Roman}

% ---- 颜色（与 figures-rendered/_style.py PALETTE='default' 同款）----
\usepackage{xcolor}
\definecolor{coBaseline}{HTML}{1f4e79}  % 深蓝 - 主体强调 / 基线
\definecolor{coImproved}{HTML}{2e7d32}  % 深绿 - 提升 / 高亮
\definecolor{coThird}   {HTML}{b8860b}  % 深金 - 第三系列
\definecolor{coFourth}  {HTML}{6a4c93}  % 深紫 - 第四系列
\definecolor{coAnno}    {HTML}{404040}  % 学术深灰 - 注记
\definecolor{coGuide}   {HTML}{808080}  % 辅助线
\definecolor{coGrid}    {HTML}{cccccc}  % 网格
\definecolor{coAccent}  {HTML}{c00000}  % 强调红 - 仅用于关键警示 / 跨节点
% 浅色填充（CMYK 转 RGB 近似）
\definecolor{flBlue}    {RGB}{226,238,247}  % 淡蓝 CMYK 30/10/0/0 ≈
\definecolor{flYellow}  {RGB}{252,242,219}  % 淡黄 CMYK 0/10/30/0 ≈
\definecolor{flGray}    {RGB}{235,235,235}  % 淡灰 CMYK 0/0/0/10 ≈
\definecolor{flGreen}   {RGB}{226,239,224}  % 浅绿 - 高亮区域

% ---- TikZ + 常用库 ----
\usepackage{tikz}
\usetikzlibrary{
    arrows.meta,
    positioning,
    calc,
    fit,
    backgrounds,
    decorations.pathreplacing,
    decorations.pathmorphing,
    shapes.geometric,
    shapes.misc,
    shapes.multipart,
    matrix,
    chains,
    shadows.blur,    % 不用阴影 (禁用), 但 fit 节点偶尔用 inner sep
    patterns
}

% ---- 全局 TikZ 默认 ----
\tikzset{
    every picture/.style={
        line cap=round,
        line join=round,
        >={Stealth[length=2.4mm, width=1.6mm]},
        font=\sffamily\small,
    },
    % ---- 节点基础类型 ----
    %
    % 主体方框 (深色边框, 白底) - 用于一般模块
    main box/.style={
        draw=coAnno,
        line width=0.7pt,
        rectangle,
        rounded corners=1.5pt,
        minimum width=20mm,
        minimum height=8mm,
        align=center,
        inner sep=2pt,
    },
    % 强调方框 (深蓝填充 alpha)
    accent box/.style={
        main box,
        fill=flBlue,
    },
    % 计算单元 / PE (淡黄填充)
    pe box/.style={
        main box,
        fill=flYellow,
    },
    % SRAM / 存储 (淡灰填充)
    mem box/.style={
        main box,
        fill=flGray,
    },
    % 控制单元 / FSM (浅绿填充)
    ctrl box/.style={
        main box,
        fill=flGreen,
    },
    % 大容器 (圆角矩形, 虚线边框, 用于 multicore_top / NUMA 节点等)
    container box/.style={
        draw=coAnno,
        line width=0.7pt,
        rounded corners=2pt,
        rectangle,
        dashed,
        inner sep=3mm,
    },
    container solid/.style={
        container box,
        solid,
    },
    %
    % ---- 箭头类型 ----
    %
    flow/.style={
        -{Stealth[length=2.4mm,width=1.6mm]},
        line width=0.7pt,
        draw=coAnno,
    },
    flow thick/.style={
        flow,
        line width=1.2pt,
    },
    flow dashed/.style={
        flow,
        dashed,
    },
    flow accent/.style={
        -{Stealth[length=2.4mm,width=1.6mm]},
        line width=1.0pt,
        draw=coAccent,
        dashed,
    },
    flow data/.style={
        -{Stealth[length=2.4mm,width=1.6mm]},
        line width=0.9pt,
        draw=coBaseline,
    },
    flow weight/.style={
        -{Stealth[length=2.4mm,width=1.6mm]},
        line width=0.9pt,
        draw=coImproved,
    },
    flow psum/.style={
        -{Stealth[length=2.4mm,width=1.6mm]},
        line width=0.9pt,
        draw=coThird,
    },
    bus/.style={
        -{Stealth[length=3mm,width=2mm]},
        line width=1.6pt,
        draw=coAnno,
    },
    handshake/.style={
        <->,
        line width=0.7pt,
        draw=coAnno,
    },
    %
    % ---- 标注 / 文字 ----
    %
    label small/.style={
        font=\sffamily\footnotesize,
        text=coAnno,
        align=center,
    },
    label tiny/.style={
        font=\sffamily\scriptsize,
        text=coAnno,
        align=center,
    },
    label bold/.style={
        font=\sffamily\bfseries\small,
        text=black,
        align=center,
    },
    formula/.style={
        font=\small,
        align=center,
    },
    % 边线上的标注 (背景白底, 不挡线)
    edge label/.style={
        font=\sffamily\scriptsize,
        text=coAnno,
        fill=white,
        inner sep=1.2pt,
        align=center,
    },
    edge label accent/.style={
        edge label,
        text=coAccent,
        font=\sffamily\bfseries\scriptsize,
    },
    %
    % ---- 图例 ----
    %
    legend box/.style={
        draw=coAnno,
        line width=0.4pt,
        fill=white,
        rounded corners=1pt,
        inner sep=2.5pt,
        font=\sffamily\scriptsize,
        text=coAnno,
        align=left,
    },
    %
    % ---- 网格 / 阵列 (用于 PE 阵列示意) ----
    %
    grid cell/.style={
        draw=coAnno,
        line width=0.3pt,
        minimum width=4mm,
        minimum height=4mm,
        inner sep=0pt,
        rectangle,
    },
    grid cell light/.style={
        grid cell,
        fill=flGray,
    },
    grid cell accent/.style={
        grid cell,
        fill=flBlue,
    },
    %
    % ---- 时序波形 (用于内嵌时序示意时, 非主要; 主要用 WaveDrom) ----
    %
    wave hi/.style={line width=1pt, draw=coBaseline},
    wave lo/.style={line width=1pt, draw=coBaseline},
}

% ---- 通用辅助宏 ----

% 章节图标题命令 (备选, 一般由论文 caption 承担, 不在 TikZ 内画)
% \newcommand{\figtitle}[1]{\node[font=\bfseries\small,text=black] {#1};}

% 端口标记: 在节点边缘画个小圆点表示端口
\newcommand{\port}[2][black]{\fill[#1] (#2) circle (0.5mm);}

% 中文/英文双行标签: \biLabel{中文}{English}
\newcommand{\biLabel}[2]{\shortstack{\sffamily #1\\\itshape\footnotesize #2}}

% =============================================================================
% 使用约定 (写 figX-Y.tex 时遵守):
%
% 1. 不在 TikZ 内画"图 X.Y"标题——由 Word caption 承担
% 2. 节点用上述预定义样式 (main box / accent box / pe box / mem box / ctrl box / container box)
%    避免每张图重复定义节点样式
% 3. 数据流统一用 flow data (蓝) / flow weight (绿) / flow psum (金) 三色编码
%    握手 / 控制流用 flow / flow dashed (灰)
%    强调 / 警示用 flow accent (红虚线), 慎用
% 4. 文字标注用 label small / label tiny / label bold / edge label
% 5. 图例统一放右下角, 用 legend box 样式
% 6. 中英混排时: 中文用 \sffamily (思源黑体), 英文数学用 \rmfamily (Times New Roman)
% =============================================================================

%   \begin{document}
%   \begin{tikzpicture}[<额外局部样式>]
%     ...
%   \end{tikzpicture}
%   \end{document}
%
% 风格基线: IEEE 期刊配色 / 白底黑字 / Times New Roman + Microsoft YaHei /
%           禁卡通 / 禁渐变 / 禁阴影 / 禁 3D / 禁 emoji
% 与 figures-rendered/_style.py (matplotlib) 配色一致, 保证全文图风格统一.
% =============================================================================

% ---- 字体（XeLaTeX 必须）----
% 中文字体：Microsoft YaHei（保持原配置，本次只调整英文）.
% 英文字体：Times New Roman, 把 \sffamily / \rmfamily 同时映射到 Times,
%           保证 TikZ 节点无论用哪个 family, 拉丁字符都是 Times New Roman.
% 数学公式（mathmode）保留 latin modern math, 与 Times 视觉差异可接受.
\usepackage{amsmath}
\usepackage{amssymb}
\usepackage{xeCJK}
\setCJKmainfont{Microsoft YaHei}
\setCJKsansfont{Microsoft YaHei}
\setmainfont{Times New Roman}
\setsansfont{Times New Roman}

% ---- 颜色（与 figures-rendered/_style.py PALETTE='default' 同款）----
\usepackage{xcolor}
\definecolor{coBaseline}{HTML}{1f4e79}  % 深蓝 - 主体强调 / 基线
\definecolor{coImproved}{HTML}{2e7d32}  % 深绿 - 提升 / 高亮
\definecolor{coThird}   {HTML}{b8860b}  % 深金 - 第三系列
\definecolor{coFourth}  {HTML}{6a4c93}  % 深紫 - 第四系列
\definecolor{coAnno}    {HTML}{404040}  % 学术深灰 - 注记
\definecolor{coGuide}   {HTML}{808080}  % 辅助线
\definecolor{coGrid}    {HTML}{cccccc}  % 网格
\definecolor{coAccent}  {HTML}{c00000}  % 强调红 - 仅用于关键警示 / 跨节点
% 浅色填充（CMYK 转 RGB 近似）
\definecolor{flBlue}    {RGB}{226,238,247}  % 淡蓝 CMYK 30/10/0/0 ≈
\definecolor{flYellow}  {RGB}{252,242,219}  % 淡黄 CMYK 0/10/30/0 ≈
\definecolor{flGray}    {RGB}{235,235,235}  % 淡灰 CMYK 0/0/0/10 ≈
\definecolor{flGreen}   {RGB}{226,239,224}  % 浅绿 - 高亮区域

% ---- TikZ + 常用库 ----
\usepackage{tikz}
\usetikzlibrary{
    arrows.meta,
    positioning,
    calc,
    fit,
    backgrounds,
    decorations.pathreplacing,
    decorations.pathmorphing,
    shapes.geometric,
    shapes.misc,
    shapes.multipart,
    matrix,
    chains,
    shadows.blur,    % 不用阴影 (禁用), 但 fit 节点偶尔用 inner sep
    patterns
}

% ---- 全局 TikZ 默认 ----
\tikzset{
    every picture/.style={
        line cap=round,
        line join=round,
        >={Stealth[length=2.4mm, width=1.6mm]},
        font=\sffamily\small,
    },
    % ---- 节点基础类型 ----
    %
    % 主体方框 (深色边框, 白底) - 用于一般模块
    main box/.style={
        draw=coAnno,
        line width=0.7pt,
        rectangle,
        rounded corners=1.5pt,
        minimum width=20mm,
        minimum height=8mm,
        align=center,
        inner sep=2pt,
    },
    % 强调方框 (深蓝填充 alpha)
    accent box/.style={
        main box,
        fill=flBlue,
    },
    % 计算单元 / PE (淡黄填充)
    pe box/.style={
        main box,
        fill=flYellow,
    },
    % SRAM / 存储 (淡灰填充)
    mem box/.style={
        main box,
        fill=flGray,
    },
    % 控制单元 / FSM (浅绿填充)
    ctrl box/.style={
        main box,
        fill=flGreen,
    },
    % 大容器 (圆角矩形, 虚线边框, 用于 multicore_top / NUMA 节点等)
    container box/.style={
        draw=coAnno,
        line width=0.7pt,
        rounded corners=2pt,
        rectangle,
        dashed,
        inner sep=3mm,
    },
    container solid/.style={
        container box,
        solid,
    },
    %
    % ---- 箭头类型 ----
    %
    flow/.style={
        -{Stealth[length=2.4mm,width=1.6mm]},
        line width=0.7pt,
        draw=coAnno,
    },
    flow thick/.style={
        flow,
        line width=1.2pt,
    },
    flow dashed/.style={
        flow,
        dashed,
    },
    flow accent/.style={
        -{Stealth[length=2.4mm,width=1.6mm]},
        line width=1.0pt,
        draw=coAccent,
        dashed,
    },
    flow data/.style={
        -{Stealth[length=2.4mm,width=1.6mm]},
        line width=0.9pt,
        draw=coBaseline,
    },
    flow weight/.style={
        -{Stealth[length=2.4mm,width=1.6mm]},
        line width=0.9pt,
        draw=coImproved,
    },
    flow psum/.style={
        -{Stealth[length=2.4mm,width=1.6mm]},
        line width=0.9pt,
        draw=coThird,
    },
    bus/.style={
        -{Stealth[length=3mm,width=2mm]},
        line width=1.6pt,
        draw=coAnno,
    },
    handshake/.style={
        <->,
        line width=0.7pt,
        draw=coAnno,
    },
    %
    % ---- 标注 / 文字 ----
    %
    label small/.style={
        font=\sffamily\footnotesize,
        text=coAnno,
        align=center,
    },
    label tiny/.style={
        font=\sffamily\scriptsize,
        text=coAnno,
        align=center,
    },
    label bold/.style={
        font=\sffamily\bfseries\small,
        text=black,
        align=center,
    },
    formula/.style={
        font=\small,
        align=center,
    },
    % 边线上的标注 (背景白底, 不挡线)
    edge label/.style={
        font=\sffamily\scriptsize,
        text=coAnno,
        fill=white,
        inner sep=1.2pt,
        align=center,
    },
    edge label accent/.style={
        edge label,
        text=coAccent,
        font=\sffamily\bfseries\scriptsize,
    },
    %
    % ---- 图例 ----
    %
    legend box/.style={
        draw=coAnno,
        line width=0.4pt,
        fill=white,
        rounded corners=1pt,
        inner sep=2.5pt,
        font=\sffamily\scriptsize,
        text=coAnno,
        align=left,
    },
    %
    % ---- 网格 / 阵列 (用于 PE 阵列示意) ----
    %
    grid cell/.style={
        draw=coAnno,
        line width=0.3pt,
        minimum width=4mm,
        minimum height=4mm,
        inner sep=0pt,
        rectangle,
    },
    grid cell light/.style={
        grid cell,
        fill=flGray,
    },
    grid cell accent/.style={
        grid cell,
        fill=flBlue,
    },
    %
    % ---- 时序波形 (用于内嵌时序示意时, 非主要; 主要用 WaveDrom) ----
    %
    wave hi/.style={line width=1pt, draw=coBaseline},
    wave lo/.style={line width=1pt, draw=coBaseline},
}

% ---- 通用辅助宏 ----

% 章节图标题命令 (备选, 一般由论文 caption 承担, 不在 TikZ 内画)
% \newcommand{\figtitle}[1]{\node[font=\bfseries\small,text=black] {#1};}

% 端口标记: 在节点边缘画个小圆点表示端口
\newcommand{\port}[2][black]{\fill[#1] (#2) circle (0.5mm);}

% 中文/英文双行标签: \biLabel{中文}{English}
\newcommand{\biLabel}[2]{\shortstack{\sffamily #1\\\itshape\footnotesize #2}}

% =============================================================================
% 使用约定 (写 figX-Y.tex 时遵守):
%
% 1. 不在 TikZ 内画"图 X.Y"标题——由 Word caption 承担
% 2. 节点用上述预定义样式 (main box / accent box / pe box / mem box / ctrl box / container box)
%    避免每张图重复定义节点样式
% 3. 数据流统一用 flow data (蓝) / flow weight (绿) / flow psum (金) 三色编码
%    握手 / 控制流用 flow / flow dashed (灰)
%    强调 / 警示用 flow accent (红虚线), 慎用
% 4. 文字标注用 label small / label tiny / label bold / edge label
% 5. 图例统一放右下角, 用 legend box 样式
% 6. 中英混排时: 中文用 \sffamily (思源黑体), 英文数学用 \rmfamily (Times New Roman)
% =============================================================================

%   \begin{document}
%   \begin{tikzpicture}[<额外局部样式>]
%     ...
%   \end{tikzpicture}
%   \end{document}
%
% 风格基线: IEEE 期刊配色 / 白底黑字 / Times New Roman + Microsoft YaHei /
%           禁卡通 / 禁渐变 / 禁阴影 / 禁 3D / 禁 emoji
% 与 figures-rendered/_style.py (matplotlib) 配色一致, 保证全文图风格统一.
% =============================================================================

% ---- 字体（XeLaTeX 必须）----
% 中文字体：Microsoft YaHei（保持原配置，本次只调整英文）.
% 英文字体：Times New Roman, 把 \sffamily / \rmfamily 同时映射到 Times,
%           保证 TikZ 节点无论用哪个 family, 拉丁字符都是 Times New Roman.
% 数学公式（mathmode）保留 latin modern math, 与 Times 视觉差异可接受.
\usepackage{amsmath}
\usepackage{amssymb}
\usepackage{xeCJK}
\setCJKmainfont{Microsoft YaHei}
\setCJKsansfont{Microsoft YaHei}
\setmainfont{Times New Roman}
\setsansfont{Times New Roman}

% ---- 颜色（与 figures-rendered/_style.py PALETTE='default' 同款）----
\usepackage{xcolor}
\definecolor{coBaseline}{HTML}{1f4e79}  % 深蓝 - 主体强调 / 基线
\definecolor{coImproved}{HTML}{2e7d32}  % 深绿 - 提升 / 高亮
\definecolor{coThird}   {HTML}{b8860b}  % 深金 - 第三系列
\definecolor{coFourth}  {HTML}{6a4c93}  % 深紫 - 第四系列
\definecolor{coAnno}    {HTML}{404040}  % 学术深灰 - 注记
\definecolor{coGuide}   {HTML}{808080}  % 辅助线
\definecolor{coGrid}    {HTML}{cccccc}  % 网格
\definecolor{coAccent}  {HTML}{c00000}  % 强调红 - 仅用于关键警示 / 跨节点
% 浅色填充（CMYK 转 RGB 近似）
\definecolor{flBlue}    {RGB}{226,238,247}  % 淡蓝 CMYK 30/10/0/0 ≈
\definecolor{flYellow}  {RGB}{252,242,219}  % 淡黄 CMYK 0/10/30/0 ≈
\definecolor{flGray}    {RGB}{235,235,235}  % 淡灰 CMYK 0/0/0/10 ≈
\definecolor{flGreen}   {RGB}{226,239,224}  % 浅绿 - 高亮区域

% ---- TikZ + 常用库 ----
\usepackage{tikz}
\usetikzlibrary{
    arrows.meta,
    positioning,
    calc,
    fit,
    backgrounds,
    decorations.pathreplacing,
    decorations.pathmorphing,
    shapes.geometric,
    shapes.misc,
    shapes.multipart,
    matrix,
    chains,
    shadows.blur,    % 不用阴影 (禁用), 但 fit 节点偶尔用 inner sep
    patterns
}

% ---- 全局 TikZ 默认 ----
\tikzset{
    every picture/.style={
        line cap=round,
        line join=round,
        >={Stealth[length=2.4mm, width=1.6mm]},
        font=\sffamily\small,
    },
    % ---- 节点基础类型 ----
    %
    % 主体方框 (深色边框, 白底) - 用于一般模块
    main box/.style={
        draw=coAnno,
        line width=0.7pt,
        rectangle,
        rounded corners=1.5pt,
        minimum width=20mm,
        minimum height=8mm,
        align=center,
        inner sep=2pt,
    },
    % 强调方框 (深蓝填充 alpha)
    accent box/.style={
        main box,
        fill=flBlue,
    },
    % 计算单元 / PE (淡黄填充)
    pe box/.style={
        main box,
        fill=flYellow,
    },
    % SRAM / 存储 (淡灰填充)
    mem box/.style={
        main box,
        fill=flGray,
    },
    % 控制单元 / FSM (浅绿填充)
    ctrl box/.style={
        main box,
        fill=flGreen,
    },
    % 大容器 (圆角矩形, 虚线边框, 用于 multicore_top / NUMA 节点等)
    container box/.style={
        draw=coAnno,
        line width=0.7pt,
        rounded corners=2pt,
        rectangle,
        dashed,
        inner sep=3mm,
    },
    container solid/.style={
        container box,
        solid,
    },
    %
    % ---- 箭头类型 ----
    %
    flow/.style={
        -{Stealth[length=2.4mm,width=1.6mm]},
        line width=0.7pt,
        draw=coAnno,
    },
    flow thick/.style={
        flow,
        line width=1.2pt,
    },
    flow dashed/.style={
        flow,
        dashed,
    },
    flow accent/.style={
        -{Stealth[length=2.4mm,width=1.6mm]},
        line width=1.0pt,
        draw=coAccent,
        dashed,
    },
    flow data/.style={
        -{Stealth[length=2.4mm,width=1.6mm]},
        line width=0.9pt,
        draw=coBaseline,
    },
    flow weight/.style={
        -{Stealth[length=2.4mm,width=1.6mm]},
        line width=0.9pt,
        draw=coImproved,
    },
    flow psum/.style={
        -{Stealth[length=2.4mm,width=1.6mm]},
        line width=0.9pt,
        draw=coThird,
    },
    bus/.style={
        -{Stealth[length=3mm,width=2mm]},
        line width=1.6pt,
        draw=coAnno,
    },
    handshake/.style={
        <->,
        line width=0.7pt,
        draw=coAnno,
    },
    %
    % ---- 标注 / 文字 ----
    %
    label small/.style={
        font=\sffamily\footnotesize,
        text=coAnno,
        align=center,
    },
    label tiny/.style={
        font=\sffamily\scriptsize,
        text=coAnno,
        align=center,
    },
    label bold/.style={
        font=\sffamily\bfseries\small,
        text=black,
        align=center,
    },
    formula/.style={
        font=\small,
        align=center,
    },
    % 边线上的标注 (背景白底, 不挡线)
    edge label/.style={
        font=\sffamily\scriptsize,
        text=coAnno,
        fill=white,
        inner sep=1.2pt,
        align=center,
    },
    edge label accent/.style={
        edge label,
        text=coAccent,
        font=\sffamily\bfseries\scriptsize,
    },
    %
    % ---- 图例 ----
    %
    legend box/.style={
        draw=coAnno,
        line width=0.4pt,
        fill=white,
        rounded corners=1pt,
        inner sep=2.5pt,
        font=\sffamily\scriptsize,
        text=coAnno,
        align=left,
    },
    %
    % ---- 网格 / 阵列 (用于 PE 阵列示意) ----
    %
    grid cell/.style={
        draw=coAnno,
        line width=0.3pt,
        minimum width=4mm,
        minimum height=4mm,
        inner sep=0pt,
        rectangle,
    },
    grid cell light/.style={
        grid cell,
        fill=flGray,
    },
    grid cell accent/.style={
        grid cell,
        fill=flBlue,
    },
    %
    % ---- 时序波形 (用于内嵌时序示意时, 非主要; 主要用 WaveDrom) ----
    %
    wave hi/.style={line width=1pt, draw=coBaseline},
    wave lo/.style={line width=1pt, draw=coBaseline},
}

% ---- 通用辅助宏 ----

% 章节图标题命令 (备选, 一般由论文 caption 承担, 不在 TikZ 内画)
% \newcommand{\figtitle}[1]{\node[font=\bfseries\small,text=black] {#1};}

% 端口标记: 在节点边缘画个小圆点表示端口
\newcommand{\port}[2][black]{\fill[#1] (#2) circle (0.5mm);}

% 中文/英文双行标签: \biLabel{中文}{English}
\newcommand{\biLabel}[2]{\shortstack{\sffamily #1\\\itshape\footnotesize #2}}

% =============================================================================
% 使用约定 (写 figX-Y.tex 时遵守):
%
% 1. 不在 TikZ 内画"图 X.Y"标题——由 Word caption 承担
% 2. 节点用上述预定义样式 (main box / accent box / pe box / mem box / ctrl box / container box)
%    避免每张图重复定义节点样式
% 3. 数据流统一用 flow data (蓝) / flow weight (绿) / flow psum (金) 三色编码
%    握手 / 控制流用 flow / flow dashed (灰)
%    强调 / 警示用 flow accent (红虚线), 慎用
% 4. 文字标注用 label small / label tiny / label bold / edge label
% 5. 图例统一放右下角, 用 legend box 样式
% 6. 中英混排时: 中文用 \sffamily (思源黑体), 英文数学用 \rmfamily (Times New Roman)
% =============================================================================


\begin{document}
\begin{tikzpicture}[
    every node/.style={font=\sffamily\small},
    dma box/.style={
        main box, fill=flBlue,
        minimum width=26mm, minimum height=12mm,
        font=\sffamily\large\bfseries,
    },
    dfe box/.style={
        main box, fill=flGreen,
        minimum width=26mm, minimum height=12mm,
        font=\sffamily\large\bfseries,
    },
    arb box/.style={
        main box, fill=flGreen,
        minimum width=36mm, minimum height=13mm,
        font=\sffamily\large\bfseries,
    },
    axidm shell/.style={
        draw=coAnno, line width=1.2pt, rounded corners=3pt,
        fill=flYellow!50, minimum width=150mm, minimum height=46mm,
    },
    sub block/.style={
        draw=coAnno, line width=0.6pt, rectangle,
        fill=white, minimum width=30mm, minimum height=11mm,
        align=center, font=\sffamily\small\bfseries,
        inner sep=2pt,
    },
    mux box/.style={
        main box, fill=flBlue,
        minimum width=140mm, minimum height=13mm,
        font=\sffamily\large\bfseries,
    },
    master box/.style={
        main box, fill=flGray,
        minimum width=60mm, minimum height=13mm,
        font=\sffamily\large\bfseries,
    },
    section title/.style={
        font=\sffamily\large\bfseries, text=coAnno, anchor=west,
    },
    big tiny/.style={
        font=\sffamily\footnotesize, text=coAnno, align=center,
    },
    big edge/.style={
        font=\sffamily\footnotesize, text=coAnno,
        fill=white, inner sep=1.5pt, align=center,
    },
    flow cmd/.style={
        -{Stealth[length=3mm,width=2mm]},
        line width=1.1pt, draw=coBaseline,
    },
    flow data big/.style={
        -{Stealth[length=3mm,width=2mm]},
        line width=1.1pt, draw=coBaseline,
    },
    flow sts/.style={
        -{Stealth[length=2.8mm,width=1.8mm]},
        line width=1.0pt, draw=coThird, dashed,
    },
    flow accent big/.style={
        -{Stealth[length=3mm,width=2mm]},
        line width=1.1pt, draw=coAccent, dashed,
    },
    bus big/.style={
        -{Stealth[length=3.5mm,width=2.4mm]},
        line width=1.8pt, draw=coAnno,
    },
]

% =============================================================================
% Layer 1: 4 DMA 控制器 + DFE (顶部)
% =============================================================================
\node[section title] at (-8.5, 7.4) {DMA 控制器};

\node[dma box] (idma) at (-6.0, 6.0) {idma\_ctrl};
\node[dma box] (wdma) at (-3.0, 6.0) {wdma\_ctrl};
\node[dma box] (rdma) at ( 0.0, 6.0) {rdma\_ctrl};
\node[dma box] (odma) at ( 4.0, 6.0) {odma\_ctrl};
\node[dfe box] (dfe)  at ( 7.5, 6.0) {DFE};

\node[big tiny, anchor=south] at (idma.north) {DDR$\to$IFB};
\node[big tiny, anchor=south] at (wdma.north) {DDR$\to$WB};
\node[big tiny, anchor=south] at (rdma.north) {DDR$\to$bias/sc};
\node[big tiny, anchor=south] at (odma.north) {OFB$\to$DDR};
\node[big tiny, anchor=south] at (dfe.north)  {desc FIFO 128};

% =============================================================================
% Layer 2: mm2s_arb (在 idma/wdma/rdma 下方居中)
% =============================================================================
\node[arb box] (arb) at (-3.0, 3.5) {mm2s\_arb};
\node[big tiny, anchor=west] at ($(arb.east)+(0.15,0)$)
    {round-robin\\3-to-1};

% 3 路 cmd → mm2s_arb
\draw[flow cmd] (idma.south) -- ($(arb.north)+(-1.5,0)$);
\draw[flow cmd] (wdma.south) -- (arb.north);
\draw[flow cmd] (rdma.south) -- ($(arb.north)+(1.5,0)$);

% =============================================================================
% Layer 3: axi_dm IP shell + 5 sub blocks
% =============================================================================
\node[section title] at (-8.5, 2.5) {axi\_dm IP \ Xilinx};

\node[axidm shell] (axidm) at (0, -0.2) {};

% sub blocks: 上排 cmd FIFO 2 个 + sts 右; 下排 data 2 个
\node[sub block] (m_cmd) at (-5.0, 1.0) {MM2S cmd FIFO};
\node[sub block] (s_cmd) at ( 0.5, 1.0) {S2MM cmd FIFO};
\node[sub block, minimum width=24mm] (sts) at ( 6.0, -0.2) {sts FIFO};
\node[sub block] (m_dat) at (-5.0, -1.4) {MM2S data};
\node[sub block] (s_dat) at ( 0.5, -1.4) {S2MM data};

% IP 内部连线: cmd FIFO → data path (同侧)
\draw[flow cmd] (m_cmd.south) -- (m_dat.north);
\draw[flow cmd] (s_cmd.south) -- (s_dat.north);

% sts FIFO 回环: 双输入走 axi_dm shell 内底部, y 错开避免重叠
% m_dat → sts.south_west: 深绕 y=-2.2
\draw[flow sts] (m_dat.east) -- (-3.1, -1.4) -- (-3.1, -2.2)
                              -- (5.6, -2.2) -- (5.6, -0.75);
% s_dat → sts.south_east: 浅绕 y=-1.8
\draw[flow sts] (s_dat.east) -- (2.4, -1.4) -- (2.4, -1.8)
                              -- (6.4, -1.8) -- (6.4, -0.75);

% mm2s_arb → MM2S cmd FIFO (右下斜)
\draw[flow cmd] (arb.south) -- (m_cmd.north)
    node[big edge, pos=0.3, right=2pt] {shared MM2S};

% odma_ctrl → S2MM cmd FIFO (直下)
\draw[flow cmd] (odma.south) -- ++(0,-2.0) -| (s_cmd.north)
    node[big edge, pos=0.15, right=2pt] {S2MM excl.};

% =============================================================================
% Layer 4: axi_m_mux (3-to-1 聚合, 底部)
% =============================================================================
\node[mux box] (mux) at (0, -4.0) {axi\_m\_mux};
\node[big tiny, anchor=west] at ($(mux.east)+(0.15,0)$) {3-to-1};

% MM2S data → mux 左
\draw[flow data big] (m_dat.south) -- ++(0,-0.6) -| ($(mux.north)+(-4.5,0)$)
    node[big edge, pos=0.92, left=2pt] {MM2S};
% S2MM data → mux 中
\draw[flow data big] (s_dat.south) -- ++(0,-0.4) -| ($(mux.north)+(0,0)$)
    node[big edge, pos=0.9, right=2pt] {S2MM};
% DFE master → mux 右 (从最右侧绕)
\draw[flow accent big] (dfe.south) -- ++(0,-9.0) -| ($(mux.north)+(4.5,0)$)
    node[big edge, pos=0.94, right=2pt, text=coAccent] {DFE};

% =============================================================================
% Layer 5: 对外 AXI4 Master → DDR
% =============================================================================
\node[master box] (mst) at (0, -5.8) {AXI4 Master $\to$ DDR};
\node[big tiny, anchor=north] at (mst.south) {128 b, via SmartConnect};

\draw[bus big] (mux.south) -- (mst.north);

% =============================================================================
% sts FIFO → DMA ctrls 完成回报 (虚线从 sts.north 直上, 拐进 odma.east)
% 标签放 sts 与 odma 之间空白处, 避开 axi_dm shell 顶部边框与 IP 标题
% =============================================================================
\draw[flow sts] (sts.north) |- (odma.east);
\node[big edge, anchor=west, text=coThird] at (6.2, 3.0)
    {sts $\to$ ctrls};

% =============================================================================
% 图例 (右下)
% =============================================================================
\node[legend box, anchor=south east, font=\sffamily\footnotesize]
    at (10.5, -7.4) {%
\begin{tabular}{@{}l@{\ \ }l@{}}
\tikz[baseline=-0.6ex]\draw[flow cmd] (0,0) -- (6mm,0); & cmd 通道 \\[1pt]
\tikz[baseline=-0.6ex]\draw[flow data big] (0,0) -- (6mm,0); & data 通道 \\[1pt]
\tikz[baseline=-0.6ex]\draw[flow sts] (0,0) -- (6mm,0); & sts 回报 \\[1pt]
\tikz[baseline=-0.6ex]\draw[flow accent big] (0,0) -- (6mm,0); & DFE 旁路 \\[1pt]
\tikz[baseline=-0.6ex]\draw[bus big] (0,0) -- (6mm,0); & 对外 AXI4 master \\
\end{tabular}%
};

\end{tikzpicture}
\end{document}
